\section{结论}\label{sec:conclusion}

本文以单源最短路为载体，把图算法的“白盒”与“黑盒”放进同一个任务、同一批图、同一套指标里做对照，得到三条结论。

\begin{enumerate}[leftmargin=1.8em]
\item \textbf{结构上，二者可以精确对应。} $k$ 层 min-sum 消息传递与同步 Bellman-Ford 的第 $k$ 轮逐点等价（定理~\ref{thm:equiv}）。
这条等价把“深度”与“边数”绑定在一起，从而给出一个构造性的深度下界：$L$ 层模型在两张 $L$ 跳计算图同构、真值不同的图上必然出错，且与训练无关（命题~\ref{prop:depth}）。
\item \textbf{白盒的鲁棒性是可证明性。} 删边单调、非树边不变、单边删除有可计算上界（命题~\ref{prop:mono}--\ref{prop:bound}），500 例验证零违反。
黑盒没有对应的命题，只能给经验分布；而在本文的实验里，它的“不敏感”其实是“没在算”。
\item \textbf{训练未必有益。} 未训练的恒等 min-sum 误差 $\sim10^{-6}$，训练 2 个参数之后反而退化到 $0.024$。
原因不是优化没做好，而是“用有限深度的监督去学一个无限迭代的算法”本身有结构性困难：损失谷底很平，学到的参数离精确解只差 0.15\%，但迭代会把这个偏差放大。
\end{enumerate}

由此得到的判断是：在最短路这类\emph{算法已知}的问题上，“可解释性 vs. 性能”的权衡曲线并不存在——白盒在准确率、鲁棒性、可解释性、计算开销四个维度上同时占优。
权衡只在算法未知时才出现；而那时，可解释性的作用不是“证明它对”，而是“发现它什么时候会错”。

\subsection*{局限与后续工作}

本文的局限已在第~\ref{sec:limits}~节逐条列出，其中最重要的两条是：黑盒基线没有调到最好（6/8 个种子塌缩），以及 L1 与 L2 的迭代预算不对称。
后续工作有三个方向：
（1）把黑盒换成权重共享的结构，让它也能迭代到收敛，从而消掉迭代预算这个不对称；
（2）把“可验证性缺口”形式化成一个可计算量（例如输出值的可达区间），而不只是一个概念；
（3）把任务从“有精确解”扩展到“有近似算法”（如旅行商），看看权衡曲线是否重新出现。

\subsection*{复现}

全部实验在纯 CPU 上完成，依赖仅为 \texttt{numpy}、\texttt{networkx}、\texttt{matplotlib} 与 \texttt{torch}，不含任何外部数据集。
仓库结构、运行命令与结果文件见 \texttt{README.md}；每个实验都有对应的验收脚本，并配有“桩实现必须失败”的负控制。
